% ID: SEL-D0034-OA-E01
% ESTADO: validada. S035, cotejada con examen y criterios.
\textbf{a)} Restamos $C$ y calculamos
\[
A^2=\begin{pmatrix}-1&-4\\2&7\end{pmatrix},\qquad
2B-C=\begin{pmatrix}5&-3&6\\6&-3&4\end{pmatrix}.
\]
Como $\det(A^2)=1$, multiplicamos por su inversa a la izquierda:
\[
X=\begin{pmatrix}7&4\\-2&-1\end{pmatrix}
\begin{pmatrix}5&-3&6\\6&-3&4\end{pmatrix}
=\boxed{\begin{pmatrix}59&-33&58\\-16&9&-16\end{pmatrix}}.
\]
Al multiplicar por $A^2$ se recupera $2B-C$.

\textbf{b)} Como $B+C$ es $2\times3$, $P$ necesita tres filas
y dos columnas para que el producto sea $2\times2$.
En $BQC^t$, las dimensiones de $B$ y $C^t$ son $2\times3$
y $3\times2$. Para que encajen ambos productos, $Q$ debe ser
$3\times3$. El resultado también es $2\times2$.
